\documentclass{article}
\usepackage[T1]{fontenc}
\usepackage{lmodern}
\usepackage{graphicx}
\usepackage{amsmath,amssymb}
\usepackage[a4paper,margin=2cm]{geometry}
\usepackage[absolute,overlay]{textpos}
\setlength{\TPHorizModule}{1mm}
\setlength{\TPVertModule}{1mm}
\usepackage[indonesian]{babel}
\usepackage{hyperref}
\setlength{\emergencystretch}{2em}
% unit: Halaman 15

\begin{document}

% === Halaman 15 (python/native) ===

Jadi, Ki merupakan penggabungan bit-bit Ci pada posisi:

14, 17,  11,  24,   1,  5,   3,  28,  15,    6,  21, 10

23, 19,  12,    4, 26,  8, 16,    7,  27,  20,  13,   2

dengan bit-bit Di pada posisi:

41, 52, 31, 37, 47, 55, 30, 40, 51,  45, 33, 48

44, 49, 39, 56, 34, 53, 46, 42, 50,  36, 29, 32

Setiap kunci internal Ki mempunyai panjang 48 bit. 

Matriks PC-2 berikut: 


14 

17 

11 

24 

1 

5 

3 

28 

15 

6 

21 

10 

23 

19 

12 

4 

26 

8 

16 

7 

27 

20 

13 

2 

41 

52 

31 

37 

47 

55 

30 

40 

51 

45 

33 

48 

44 

49 

39 

56 

34 

53 

46 

42 

50 

36 

29 

32 


15

Kunci eksternal

Permutasi

PC-1

C0

D0

Left Shift

Left Shift

C1

D1

Left Shift

Left Shift

Permutasi

PC-2

K1

Cj

Dj

Permutasi

PC-2

Kj









Left Shift

Left Shift

C16

D16

Permutasi

PC-2

K16

\end{document}
